\documentclass[../main.tex]{subfiles}
\begin{document}

\textit{Dựa trên các giới hạn đã phân tích ở Chương 2, chương này trình bày kiến trúc phương pháp luận cốt lõi của CogMem: một hệ thống bộ nhớ hội thoại dài hạn dựa trên đồ thị dị thể, retain pipeline nhiều bước và recall pipeline có định tuyến thích ứng.}

\section{Phương pháp đề xuất}

\subsection{Tổng quan kiến trúc hệ thống CogMem}

CogMem được thiết kế như một kiến trúc gồm ba pipeline liên tiếp:

\begin{enumerate}
    \item \textbf{Retain pipeline}: chuẩn hóa hội thoại, trích xuất fact, sinh embedding, lưu node và tạo liên kết.
    \item \textbf{Recall pipeline}: chạy song song nhiều kênh truy hồi, phân tích loại truy vấn, rồi dung hợp kết quả theo profile phù hợp.
    \item \textbf{Reflect / Synthesis pipeline}: tổng hợp evidence được recall để sinh câu trả lời bám dữ kiện.
\end{enumerate}

\begin{figure}[H]
\centering
\begin{tikzpicture}[x=0.01\textwidth,y=0.01\textwidth]
\node[anchor=south west,inner sep=0] at (0,0) {\includegraphics[width=\textwidth]{Images/cogmem_pipeline_overview.png}};
\node[reportlabel] at (10,51) {Conversation};
\node[reportlabel] at (29,51) {Retain};
\node[reportlabel] at (50,51) {Memory Graph};
\node[reportlabel] at (71,51) {Recall Channels};
\node[reportlabel] at (90,51) {Grounded Answer};
\node[reportnote] at (49,7) {new turns update memory};
\end{tikzpicture}
\caption{Tổng quan pipeline CogMem: hội thoại được retain thành memory graph, sau đó recall đa kênh và synthesis dựa trên evidence để trả lời truy vấn.}
\label{fig:ch3_pipeline_overview}
\end{figure}

Trong phạm vi luận án này, \textbf{BFS graph retrieval với SUM spreading activation và ba cycle guards} được xem là đường triển khai chuẩn của CogMem. Một số retriever khác được thêm vào ở các giai đoạn phát triển sau nhằm phục vụ nghiên cứu và ablation, nhưng không thay đổi lõi phương pháp luận được trình bày ở đây.

\subsubsection{Phạm vi và giả định triển khai}

Để giữ phần methodology trung thực, cần phân biệt rõ giữa \textit{ý tưởng kiến trúc} và \textit{đường triển khai được đánh giá trong luận án}. CogMem không giả định rằng mọi fact type đều luôn cần thiết trong mọi dữ liệu, và cũng không giả định rằng LLM extraction luôn phân loại đúng. Thay vào đó, hệ thống được thiết kế để cung cấp một cấu trúc biểu diễn giàu hơn bộ nhớ phẳng, rồi kiểm tra thực nghiệm xem cấu trúc đó giúp ích trong những loại câu hỏi nào.

Các giả định triển khai chính gồm:

\begin{itemize}
    \item Fact extraction vẫn phụ thuộc vào LLM, vì vậy typed node có thể bị thiếu, bị gộp hoặc bị re-typed thành fact type khác.
    \item Graph recall không thay thế semantic/BM25 retrieval; trong evaluation cuối, multi-channel recall mới là cấu hình chính.
    \item \textit{raw\_snippet} được giữ như bằng chứng nguyên văn, nhưng generation có dùng trực tiếp snippet hay không phụ thuộc vào policy của từng run.
    \item Các fact type như \textit{habit}, \textit{intention} và \textit{action\_effect} nên được hiểu là các slot biểu diễn có điều kiện, không phải các module luôn tạo uplift độc lập.
\end{itemize}

Do đó, đóng góp của chương này là mô tả một methodology có thể kiểm chứng: retain tạo typed memory graph, recall dùng nhiều kênh để lấy evidence, và reflect sinh câu trả lời dựa trên evidence đó. Các claims mạnh hơn về tác động của từng fact type phải được rút ra từ Chương 5, không được mặc định từ thiết kế.

\subsubsection{Retain pipeline}

Retain pipeline hiện thực hóa bộ nhớ dài hạn theo chuỗi bước sau:

\begin{enumerate}
    \item \textbf{Chunking theo hội thoại}: hội thoại được chia thành các đoạn phù hợp cho extraction.
    \item \textbf{Pass 1 extraction}: LLM đọc ngữ cảnh mixed-speaker để trích xuất đủ sáu fact type, đặc biệt là \textit{world} và \textit{action\_effect}.
    \item \textbf{Pass 2 extraction}: một pass persona-focused đọc lại ngữ cảnh người dùng để tinh chỉnh các fact cá nhân như \textit{experience}, \textit{habit}, \textit{intention} và \textit{opinion}.
    \item \textbf{Dedup và chuẩn hóa metadata}: hợp nhất kết quả giữa các pass, chuẩn hóa status, quan hệ chuyển pha và trường metadata chuyên biệt.
    \item \textbf{Embedding + storage}: mỗi fact được gắn embedding và lưu vào đồ thị bộ nhớ cùng \textit{raw\_snippet}.
    \item \textbf{Link creation}: tạo các liên kết entity, temporal, semantic, causal, s\_r\_link, a\_o\_causal và transition.
\end{enumerate}

Điểm quan trọng là retain pipeline không chỉ lưu một câu tóm tắt ngữ nghĩa, mà còn bảo toàn văn bản nguồn để phục vụ giai đoạn sinh câu trả lời về sau.

\subsubsection{Recall pipeline}

Recall pipeline của CogMem chạy song song bốn kênh:
\begin{itemize}
    \item semantic retrieval,
    \item BM25 retrieval,
    \item graph retrieval,
    \item temporal retrieval.
\end{itemize}

Trước khi fusion, hệ thống phân tích truy vấn để xác định \textit{query type} và \textit{temporal constraint} nếu có. Kết quả phân tích này quyết định:
\begin{itemize}
    \item fact types nào cần ưu tiên,
    \item profile trọng số RRF nào cần dùng,
    \item và có cần hậu xử lý đặc biệt cho bằng chứng prospective hoặc causal hay không.
\end{itemize}

\subsubsection{Reflect / Synthesis pipeline}

Sau khi recall hoàn tất, hệ thống chọn top evidence để đưa vào prompt sinh câu trả lời. Trong luồng này, \textit{raw\_snippet} đóng vai trò là lớp bằng chứng nguyên văn. Tùy cấu hình của giai đoạn sinh, snippet có thể được đưa thẳng vào prompt hoặc chỉ được giữ như nguồn tham chiếu nội bộ; tuy nhiên, ở mức kiến trúc, trường \textit{raw\_snippet} luôn được bảo toàn end-to-end để tránh mất dữ kiện.

\subsection{Đóng góp 1: Đồ thị Bộ nhớ Kế thừa Nhận thức (Cognitively-Grounded Memory Graph)}

\subsubsection{Nền tảng tâm lý}

Động lực của CogMem đến từ quan sát rằng bộ nhớ người không phải một kho đồng nhất. Các nhóm ký ức như semantic memory, episodic memory, habit memory, prospective memory hay action-effect binding có vai trò khác nhau trong suy luận và hành vi. Vì vậy, việc ép mọi dữ liệu hội thoại vào cùng một dạng fact sẽ làm suy yếu tính phân biệt về chức năng.

CogMem chuyển trực giác này thành một memory graph dị thể gồm nhiều fact type chuyên biệt. Mỗi network không chỉ khác nhau về nhãn, mà còn khác nhau ở metadata, loại liên kết, và chiến lược truy hồi được ưu tiên khi đối mặt với từng loại câu hỏi.

\begin{figure}[H]
\centering
\begin{tikzpicture}[x=0.01\textwidth,y=0.01\textwidth]
\node[anchor=south west,inner sep=0] at (0,0) {\includegraphics[width=\textwidth]{Images/cogmem_memory_graph.png}};
\node[reportlabel] at (27,50) {World};
\node[reportlabel] at (72,50) {Experience};
\node[reportlabel] at (18,33) {Opinion};
\node[reportlabel] at (82,33) {Habit};
\node[reportlabel] at (35,14) {Intention};
\node[reportlabel] at (65,14) {Action-Effect};
\node[reportnote] at (50,31) {typed edges + shared entities};
\end{tikzpicture}
\caption{Sáu mạng bộ nhớ của CogMem}
\label{fig:ch3_six_networks}
\end{figure}

\subsubsection{Sáu mạng bộ nhớ}

CogMem hiện dùng sáu fact type cốt lõi:

\begin{enumerate}
    \item \textbf{World}: kiến thức khách quan hoặc tương đối ổn định về thế giới.
    \item \textbf{Experience}: sự kiện cá nhân đã xảy ra, thường gắn với thời gian hoặc bối cảnh cụ thể.
    \item \textbf{Opinion}: nhận định, sở thích hoặc đánh giá chủ quan.
    \item \textbf{Habit}: mẫu hành vi lặp lại mang tính stimulus-response.
    \item \textbf{Intention}: kế hoạch hoặc mục tiêu hướng tới tương lai, có trạng thái \textit{planning / fulfilled / abandoned}.
    \item \textbf{Action-Effect}: bộ ba nhân quả \textit{precondition -- action -- outcome}.
\end{enumerate}

\begin{table}[H]
\centering
\caption{Các fact type cốt lõi và metadata chuyên biệt}
\label{tab:node_schemas}
\begin{tabular}{|p{2.7cm}|p{5.0cm}|p{4.8cm}|}
\hline
\textbf{Fact type} & \textbf{Vai trò biểu diễn} & \textbf{Metadata / tín hiệu chuyên biệt} \\
\hline
world & Kiến thức khách quan, mô tả ổn định & text, entities, raw\_snippet \\
\hline
experience & Sự kiện cá nhân đã xảy ra & occurred\_start, occurred\_end, mentioned\_at \\
\hline
opinion & Đánh giá hoặc sở thích chủ quan & confidence \\
\hline
habit & Mẫu hành vi lặp lại & liên kết s\_r\_link đến các experience liên quan \\
\hline
intention & Kế hoạch và mục tiêu tương lai & intention\_status; có thể mang thêm mô tả ``why'' hoặc tín hiệu thời gian từ extraction \\
\hline
action\_effect & Quy luật nhân quả ở cấp hành động & precondition, action, outcome, confidence, devalue\_sensitive \\
\hline
\end{tabular}
\end{table}

Việc tách riêng \textit{habit}, \textit{intention} và \textit{action\_effect} không nên được hiểu như một khẳng định rằng ba mạng này luôn tạo ra cải thiện đồng đều trên mọi tập dữ liệu. Thay vào đó, chúng cung cấp các ``representational slots'' mới cho những tín hiệu mà \textit{world} hoặc \textit{experience} thường mô tả không đầy đủ, đặc biệt ở các câu hỏi preference, prospective và causal. Trong các thí nghiệm hiện tại, \textit{habit} là nhóm có bằng chứng uplift yếu nhất: nó hợp lý về mặt biểu diễn, nhưng benchmark chưa có đủ câu hỏi thật sự cần mô hình hóa routine lặp lại để chứng minh đóng góp độc lập.

\subsubsection{Phân biệt Habit vs. Action-Effect}

CogMem tách \textit{habit} và \textit{action\_effect} vì hai loại này phản ánh hai cơ chế nhận thức khác nhau:

\begin{itemize}
    \item \textbf{Habit}: phản ứng lặp lại, mang tính S-R, hình thành từ tần suất và ngữ cảnh quen thuộc.
    \item \textbf{Action-Effect}: quy luật nhân quả có chủ đích, nơi một hành động được chọn vì tạo ra một kết quả quan sát được.
\end{itemize}

Trong thực tế hội thoại, hai loại tín hiệu này dễ bị trộn lẫn nếu chỉ dùng fact phẳng. Việc tách riêng giúp truy hồi preference hoặc causal evidence rõ hơn, đồng thời mở đường cho các thí nghiệm ablation theo từng cấu phần.

\begin{figure}[H]
\centering
\resizebox{0.92\textwidth}{!}{%
\begin{tikzpicture}
\node[reportboxgreen, minimum width=3.0cm] (stim) at (0,1.25) {Stimulus\\context};
\node[reportboxgreen, minimum width=3.0cm] (routine) at (4.3,1.25) {Repeated\\response};
\node[reportboxgold, minimum width=3.0cm] (habit) at (8.6,1.25) {Habit node\\S-R link};
\draw[reportarrow] (stim) -- (routine);
\draw[reportarrow] (routine) -- (habit);

\node[reportboxred, minimum width=3.0cm] (pre) at (0,-1.25) {Precondition\\P};
\node[reportboxred, minimum width=3.0cm] (act) at (4.3,-1.25) {Chosen action\\A};
\node[reportboxred, minimum width=3.0cm] (out) at (8.6,-1.25) {Observed\\outcome O};
\node[reportboxcyan, minimum width=3.2cm] (ae) at (13.0,-1.25) {Action-effect\\P--A--O};
\draw[reportarrow] (pre) -- (act);
\draw[reportarrow] (act) -- (out);
\draw[reportarrow] (out) -- (ae);

\node[font=\bfseries] at (-1.8,1.25) {Habit};
\node[font=\bfseries] at (-1.8,-1.25) {Action-effect};
\node[font=\scriptsize, align=center] at (4.3,0.35) {generalizes repeated behavior};
\node[font=\scriptsize, align=center] at (6.5,-2.05) {stores intentional causal intervention};
\end{tikzpicture}%
}
\caption{Phân biệt Habit và Action-Effect}
\label{fig:ch3_sr_vs_ao}
\end{figure}

\subsubsection{Bảy loại liên kết (Edge Types)}

Trên đồ thị CogMem, thông tin được kết nối bởi bảy nhóm liên kết chính:

\begin{enumerate}
    \item \textbf{entity}: nối các node cùng nhắc tới một thực thể.
    \item \textbf{temporal}: mô tả quan hệ trước-sau giữa các sự kiện.
    \item \textbf{semantic}: nối các node có tương đồng ngữ nghĩa cao.
    \item \textbf{causal}: quan hệ nhân quả hoặc ảnh hưởng giữa các fact.
    \item \textbf{s\_r\_link}: liên kết hỗ trợ habit bằng các experience cùng mẫu kích thích-phản ứng.
    \item \textbf{a\_o\_causal}: liên kết từ action\_effect tới các fact quan sát được làm bằng chứng outcome.
    \item \textbf{transition}: liên kết chuyển pha trạng thái của một biểu diễn theo thời gian.
\end{enumerate}

Riêng \textit{transition} là một edge type duy nhất nhưng mang \textit{transition\_type} bên trong. Trong implementation hiện tại, các typed transition cốt lõi gồm:
\begin{itemize}
    \item \texttt{fulfilled\_by}: intention $\rightarrow$ experience,
    \item \texttt{triggered}: experience $\rightarrow$ intention,
    \item \texttt{enabled\_by}: intention $\rightarrow$ intention,
    \item \texttt{revised\_to}: opinion $\rightarrow$ opinion,
    \item \texttt{contradicted\_by}: world $\rightarrow$ world.
\end{itemize}

\textbf{Lưu ý:} \textit{abandoned} là một \textit{intention\_status}, không phải một transition edge riêng.

\begin{figure}[H]
\centering
\resizebox{0.92\textwidth}{!}{%
\begin{tikzpicture}
\node[reportboxgold, minimum width=3.1cm, minimum height=1.0cm] (center) at (0,0) {Memory\\graph};
\node[reportbox, minimum width=3.0cm] (entity) at (-4.8,1.9) {\textbf{entity}\\[-0.1em]\scriptsize shared people, places, tools};
\node[reportbox, minimum width=3.0cm] (temporal) at (0,2.35) {\textbf{temporal}\\[-0.1em]\scriptsize before/after/session order};
\node[reportbox, minimum width=3.0cm] (semantic) at (4.8,1.9) {\textbf{semantic}\\[-0.1em]\scriptsize meaning similarity};
\node[reportbox, minimum width=3.0cm] (causal) at (-5.2,-0.25) {\textbf{causal}\\[-0.1em]\scriptsize why and because};
\node[reportbox, minimum width=3.0cm] (sr) at (5.2,-0.25) {\textbf{s\_r\_link}\\[-0.1em]\scriptsize habit evidence};
\node[reportbox, minimum width=3.0cm] (ao) at (-3.0,-2.25) {\textbf{a\_o\_causal}\\[-0.1em]\scriptsize action-outcome support};
\node[reportbox, minimum width=3.0cm] (transition) at (3.0,-2.25) {\textbf{transition}\\[-0.1em]\scriptsize state lifecycle};
\foreach \n in {entity,temporal,semantic,causal,sr,ao,transition}{\draw[reportarrow] (center) -- (\n);}
\end{tikzpicture}%
}
\caption{Bảy loại liên kết trong CogMem}
\label{fig:ch3_edge_types}
\end{figure}

\subsection{Đóng góp 2: Siêu dữ liệu Node phi tổn hao (Lossless Node Metadata) và Raw Snippet}

\subsubsection{Vấn đề: nén dữ liệu có tổn hao}

Các hệ thống bộ nhớ dựa trên extraction thường phải nén hội thoại dài thành một số fact ngắn hơn để có thể lập chỉ mục và truy hồi hiệu quả. Tuy nhiên, phép nén này không đảo ngược được. Những dữ kiện như con số chính xác, cụm từ nguyên văn, trạng thái cảm xúc hoặc các tín hiệu ``nhỏ nhưng quyết định'' thường biến mất khỏi bản tóm tắt.

Ví dụ:
\begin{lstlisting}
Original conversation:
"I just quit my job at VCCorp today.
I'm really sad because my old team was
like family. DI offered me 40% more pay.
I start there on April 1st."

Compressed fact:
"User switched from VCCorp to DI in
April 2024."
\end{lstlisting}

Fact nén trên đủ cho semantic retrieval cơ bản, nhưng không còn bảo toàn những chi tiết như tỷ lệ tăng lương, trạng thái cảm xúc hay chính xác ngày khởi đầu.

\subsubsection{Giải pháp: 2-layer node schema}

CogMem giải quyết vấn đề này bằng cách giữ đồng thời hai lớp biểu diễn:
\begin{itemize}
    \item \textbf{fact text / embedding}: lớp nén dùng cho retrieval.
    \item \textbf{raw\_snippet}: lớp nguyên văn dùng cho grounding khi cần sinh câu trả lời chi tiết.
\end{itemize}

Minh họa:
\begin{lstlisting}
{
  "fact_type": "experience",
  "text": "User switched from VCCorp to DI
           in April 2024",
  "embedding": [...],
  "raw_snippet": "I just quit my job at
                  VCCorp today... DI offered
                  me 40% more pay...",
  "metadata": {
    "raw_snippet_present": true
  }
}
\end{lstlisting}

Thiết kế này cho phép CogMem vừa truy hồi nhanh trên representation đã nén, vừa giữ được bằng chứng nguyên văn khi giai đoạn synthesis cần chi tiết chính xác. Ở implementation hiện tại, snippet có thể được đưa trực tiếp vào prompt generate khi policy của evaluation hoặc runtime bật chế độ sử dụng verbatim evidence; nhưng ngay cả khi policy đó thay đổi, \textit{raw\_snippet} vẫn được bảo toàn như một thành phần lõi của dữ liệu.

\begin{figure}[H]
\centering
\resizebox{0.92\textwidth}{!}{%
\begin{tikzpicture}
\node[reportboxcyan, minimum width=4.0cm, minimum height=1.25cm] (chat) at (0,0)
{Conversation chunk\\\scriptsize original user wording};
\node[reportboxgreen, minimum width=4.6cm, minimum height=1.45cm] (fact) at (5.2,1.15)
{Layer 1: compressed fact\\\scriptsize text, type, entities, date, embedding};
\node[reportboxgold, minimum width=4.6cm, minimum height=1.45cm] (snippet) at (5.2,-1.15)
{Layer 2: raw snippet\\\scriptsize verbatim detail, source, document id};
\node[reportbox, minimum width=3.9cm, minimum height=1.15cm] (recall) at (10.6,1.15)
{Retrieval\\\scriptsize semantic/BM25/graph/time};
\node[reportboxred, minimum width=3.9cm, minimum height=1.15cm] (answer) at (10.6,-1.15)
{Grounded answer\\\scriptsize exact numbers and names};
\draw[reportarrow] (chat) -- (fact);
\draw[reportarrow] (chat) -- (snippet);
\draw[reportarrow] (fact) -- (recall);
\draw[reportarrow] (snippet) -- (answer);
\draw[reportarrow] (recall) -- (answer);
\end{tikzpicture}%
}
\caption{2-Layer Node Schema}
\label{fig:ch3_two_layer_schema}
\end{figure}

\subsection{Đóng góp 3: Episodic Buffer với SUM Spreading Activation + Cycle Guards}

\subsubsection{Vấn đề: MAX chỉ giữ đường mạnh nhất}

Một hạn chế của graph retrieval kiểu truyền thống là chỉ cho phép một nguồn bằng chứng ``mạnh nhất'' quyết định activation ở node đích. Dạng công thức này phù hợp với trường hợp evidence tập trung, nhưng dễ làm mất tín hiệu khi câu trả lời phụ thuộc vào nhiều đường yếu bổ sung cho nhau.

\begin{equation}
A(v, t+1) = \max_{u \in N(v)}
\left[
A(u, t) \cdot w(u,v) \cdot \delta \cdot \mu(\ell)
\right]
\end{equation}

Trong hội thoại dài hạn, đây là vấn đề thực sự vì thông tin trả lời thường bị chia nhỏ qua nhiều phiên, nhiều node và nhiều loại quan hệ.

\subsubsection{Giải pháp: SUM thay cho MAX}

CogMem dùng tư duy của \textit{episodic buffer}: activation ở một node nên có thể cộng dồn từ nhiều nguồn bằng chứng liên quan. Trên BFS graph retriever mặc định, CogMem dùng phép SUM với clipping để chặn bùng nổ tín hiệu:

\begin{equation}
\begin{aligned}
A(v, t+1) = \mathrm{clip}\Bigg(
& A(v,t) + \delta \cdot \sum_{u \in N(v)}
\Big[
A(u,t) \cdot w(u,v) \\
& \cdot \mu(\mathrm{edge}) \cdot
\mathrm{refractory}(u)
\Big],
\; A_{\max}
\Bigg)
\end{aligned}
\end{equation}

\begin{table}[H]
\centering
\caption{So sánh HINDSIGHT (MAX) và CogMem (SUM)}
\label{tab:max_vs_sum}
\begin{tabular}{|p{3.2cm}|p{4.2cm}|p{4.2cm}|}
\hline
\textbf{Khía cạnh} & \textbf{HINDSIGHT (MAX)} & \textbf{CogMem (SUM)} \\
\hline
Bằng chứng đa đường & chỉ giữ đường tốt nhất & cộng dồn tín hiệu từ nhiều đường \\
\hline
Multi-hop reasoning & dễ rơi khi một mắt xích yếu & bền hơn với evidence phân tán \\
\hline
Yêu cầu entry point & cần trúng node rất mạnh ngay từ đầu & cho phép nhiều tín hiệu vừa phải bổ sung nhau \\
\hline
Rủi ro kỹ thuật & ổn định hơn nhưng bảo thủ & mạnh hơn nhưng cần cycle guards \\
\hline
\end{tabular}
\end{table}

Để kiểm tra đóng góp này một cách công bằng, phần thực nghiệm bổ sung trong Chương 5 dùng một đối chứng hẹp: cùng retained memory banks, cùng graph-only recall, cùng tập fact type, cùng bỏ reranker, chỉ thay reducer của BFS graph activation từ SUM sang MAX. Vì vậy, khác biệt đo được trong thí nghiệm này không phải do semantic retrieval, BM25, temporal retrieval hay generation, mà chủ yếu phản ánh cách graph channel xếp hạng evidence khi nhiều đường yếu cùng hội tụ về một node trả lời.

\subsubsection{Ba cơ chế bảo vệ chu trình (Cycle Guards)}

SUM activation trên đồ thị có chu trình cần cơ chế ổn định hóa. CogMem dùng ba guard bổ sung cho nhau:

\begin{enumerate}
    \item \textbf{Local Refractory Period}: một node vừa fire sẽ không được tái kích hoạt ngay ở bước kế tiếp.
    \begin{equation}
    \mathrm{refractory}(u) =
    \begin{cases}
    0, & \text{nếu } u \text{ vừa kích hoạt ở bước trước} \\
    1, & \text{ngược lại}
    \end{cases}
    \end{equation}

    \item \textbf{Global Firing Quota}: mỗi node chỉ được nhận số lần kích hoạt hữu hạn trong một đợt propagation, nhờ đó chặn chu trình dài.

    \item \textbf{Saturation Threshold}: giá trị activation bị chặn trên để tránh score explosion và duy trì đa dạng kết quả retrieval.
\end{enumerate}

Ba guard này phản ánh lập luận nhận thức rằng hệ thống truy hồi không nên bị ``kẹt'' trong một vòng reverberation vô hạn chỉ vì đồ thị có nhiều vòng lặp cục bộ.

\begin{figure}[H]
\centering
\resizebox{0.92\textwidth}{!}{%
\begin{tikzpicture}
\node[reportboxcyan, minimum width=3.2cm] (signal) at (0,0) {Incoming\\activation};
\node[reportboxgreen, minimum width=3.2cm] (refr) at (4.0,1.55) {Refractory\\\scriptsize block immediate refire};
\node[reportboxgold, minimum width=3.2cm] (quota) at (4.0,0) {Firing quota\\\scriptsize finite activations};
\node[reportboxred, minimum width=3.2cm] (sat) at (4.0,-1.55) {Saturation\\\scriptsize cap score};
\node[reportbox, minimum width=3.5cm] (rank) at (8.4,0) {Stable ranked\\evidence list};
\draw[reportarrow] (signal) -- (refr);
\draw[reportarrow] (signal) -- (quota);
\draw[reportarrow] (signal) -- (sat);
\draw[reportarrow] (refr) -- (rank);
\draw[reportarrow] (quota) -- (rank);
\draw[reportarrow] (sat) -- (rank);
\draw[reportarrow, dashed] (quota.south) .. controls (2.1,-2.6) and (2.1,2.6) .. node[left, font=\scriptsize]{cycle guard} (refr.north);
\end{tikzpicture}%
}
\caption{Ba cơ chế bảo vệ chu trình}
\label{fig:ch3_cycle_guards}
\end{figure}

\subsection{Đóng góp 4: Adaptive Query Routing}

\subsubsection{Vấn đề: equal-weight fusion không phản ánh ý định truy vấn}

Nếu mọi câu hỏi đều dùng cùng một profile fusion, hệ thống sẽ đối xử như nhau với các truy vấn semantic, temporal, causal hay prospective. Điều đó mâu thuẫn với trực giác nhận thức: câu hỏi ``khi nào'' nên ưu tiên temporal evidence nhiều hơn, trong khi câu hỏi ``tại sao'' cần graph evidence và causal links nhiều hơn.

\subsubsection{Giải pháp: định tuyến và dung hợp theo query type}

CogMem dùng một tầng \textit{query analysis} trước fusion. Bộ phân tích này kết hợp các lexical rules với temporal extraction để gán truy vấn vào một trong sáu loại:
\begin{itemize}
    \item semantic,
    \item temporal,
    \item causal,
    \item prospective,
    \item preference,
    \item multi\_hop.
\end{itemize}

Mỗi loại truy vấn được gắn một profile trọng số riêng cho bốn channel. Các hệ số dưới đây là \textbf{trọng số tương đối} trong weighted RRF, không phải xác suất chuẩn hóa bắt buộc có tổng bằng 1:

\begin{equation}
\mathrm{RRF}_{\mathrm{adaptive}}(d, q) =
\sum_i
w_i(q) \cdot \frac{1}{60 + \mathrm{rank}_i(d)}
\end{equation}

\begin{table}[H]
\centering
\caption{Adaptive RRF weights theo query type trong implementation hiện tại}
\label{tab:adaptive_rrf}
\resizebox{\linewidth}{!}{%
\begin{tabular}{|l|c|c|c|c|p{4.0cm}|}
\hline
\textbf{Query Type} & $w_\text{sem}$ & $w_\text{bm25}$ & $w_\text{graph}$ & $w_\text{temp}$ & \textbf{Ghi chú} \\
\hline
semantic & 1.0 & 1.0 & 1.0 & 1.0 & baseline profile \\
\hline
temporal & 0.8 & 0.6 & 0.8 & 2.2 & ưu tiên temporal filtering và temporal evidence \\
\hline
causal & 0.8 & 0.7 & 2.4 & 1.0 & ưu tiên graph channel và action-effect evidence \\
\hline
prospective & 0.9 & 0.8 & 2.0 & 1.4 & ưu tiên planning intentions và transition evidence \\
\hline
preference & 1.0 & 1.2 & 1.4 & 0.5 & ưu tiên graph / habit-related evidence \\
\hline
multi\_hop & 0.9 & 0.7 & 2.6 & 0.8 & ưu tiên graph traversal đa bước \\
\hline
\end{tabular}% 
}
\end{table}

Ngoài profile trọng số, implementation hiện tại còn có hai hậu xử lý đáng chú ý:
\begin{itemize}
    \item \textbf{Prospective filtering}: nếu truy vấn được phân loại là \textit{prospective}, hệ thống ưu tiên \textit{intention} nodes có \textit{intention\_status=planning}. Nếu một ablation loại bỏ hẳn intention network, recall sẽ fallback sang các fact type còn lại để tránh chiến thắng giả do router thiên vị.
    \item \textbf{Evidence priority}: sau weighted RRF, causal queries được boost thêm khi evidence đến từ \textit{action\_effect} hoặc graph-backed chains; prospective queries được boost thêm khi evidence là planning intention có graph/temporal backing.
\end{itemize}

\begin{figure}[H]
\centering
\resizebox{0.92\textwidth}{!}{%
\begin{tikzpicture}
\node[reportboxcyan, minimum width=3.2cm] (query) at (0,0) {User query};
\node[reportboxgold, minimum width=3.4cm] (analyzer) at (4.0,0) {Query analyzer\\\scriptsize lexical + temporal cues};
\node[reportbox, minimum width=2.25cm] (sem) at (8.2,1.9) {semantic};
\node[reportbox, minimum width=2.25cm] (temp) at (10.9,1.9) {temporal};
\node[reportbox, minimum width=2.25cm] (causal) at (8.2,0.65) {causal};
\node[reportbox, minimum width=2.25cm] (pros) at (10.9,0.65) {prospective};
\node[reportbox, minimum width=2.25cm] (pref) at (8.2,-0.6) {preference};
\node[reportbox, minimum width=2.25cm] (multi) at (10.9,-0.6) {multi-hop};
\node[reportboxgreen, minimum width=4.0cm] (rrf) at (9.55,-2.25) {Weighted RRF\\\scriptsize semantic + BM25 + graph + time};
\draw[reportarrow] (query) -- (analyzer);
\foreach \n in {sem,temp,causal,pros,pref,multi}{\draw[reportarrow] (analyzer) -- (\n);}
\foreach \n in {sem,temp,causal,pros,pref,multi}{\draw[reportarrow] (\n) -- (rrf);}
\end{tikzpicture}%
}
\caption{Adaptive Query Routing}
\label{fig:ch3_adaptive_routing}
\end{figure}


\vspace{1em}
\noindent\textit{Sau khi mô tả kiến trúc và bốn đóng góp cốt lõi của CogMem, Chương 4 sẽ chuyển sang phân tích cơ sở khoa học nhận thức của từng quyết định thiết kế.}

\end{document}
